\chapter*{Abstract}
\addcontentsline{toc}{chapter}{Abstract}

This report presents a determinant-based workflow for predicting antimicrobial resistance in \textit{Salmonella enterica} from bacterial genomic data. The main objective is to determine whether curated resistance determinants, including acquired resistance genes and resistance-associated mutations, can be used to predict binary resistant/susceptible phenotypes for individual antibiotics. Instead of relying on genome-wide feature spaces that are harder to interpret biologically, the study focuses on a curated determinant space derived from AMRFinderPlus and MicroBIGG-E.

The workflow has three main parts: cleaning phenotype labels from AST data, cleaning and normalizing genotype records, and building one model-ready table for each antibiotic. To study how much biological knowledge should be used before learning starts, the report compares three connection scopes: \textit{strict}, \textit{broad}, and \textit{all}. XGBoost models are then trained for each antibiotic with five-fold cross-validation. In addition, two simple rule-based baselines, \textit{strict rule-based} and \textit{broad rule-based}, are used as direct references without machine learning.

The results show that \textit{all ML} gives the best overall predictive performance, while \textit{broad ML} gives the best balance between feature coverage and biological interpretability. The rule-based baselines also show that, for some antibiotics, known genes and mutations are already enough for very accurate prediction. For other antibiotics, however, this direct gene-to-label rule is too simple and produces many false positives or false negatives. In those cases, machine learning helps by combining several signals at the same time.

Overall, the study shows that AMR prediction from curated genes and mutations is a suitable bioinformatics approach for this problem. It also shows that the difference between \textit{strict}, \textit{broad}, and \textit{all} is important: these settings represent different trade-offs between biological simplicity, ease of interpretation, and predictive accuracy.
